\documentclass[a4paper]{article}
% Shared preamble for the labs (compile from Labs/ with XeLaTeX or tectonic).
% Same fonts and header as the Series sheets, plus colored boxes and tables.

\usepackage{enumitem}
\usepackage[margin=2.2cm]{geometry}
\usepackage{amsmath,amsfonts,amssymb}
\usepackage[charter]{mathdesign}
\usepackage{array}
\usepackage{colortbl}
\usepackage{fancyhdr}
\usepackage[skins,breakable]{tcolorbox}
\usepackage{polyglossia}

%------------------------------------------------
% Languages and fonts
%------------------------------------------------
\setdefaultlanguage[calendar=gregorian, numerals=maghrib]{arabic}
\setotherlanguage{english}

\IfFontExistsTF{Amiri}{%
    \newfontfamily\arabicfont[Script=Arabic, Scale=1.15]{Amiri}
}{%
    \newfontfamily\arabicfont[
        Script=Arabic, Scale=1.15,
        Path=../fonts/, Extension=.ttf,
        UprightFont=*-Regular, BoldFont=*-Bold,
        ItalicFont=*-Italic, BoldItalicFont=*-BoldItalic
    ]{Amiri}
}

\IfFontExistsTF{AlBattar}{%
    \newfontfamily\arabicfontsf[Script=Arabic, Scale=1.2]{AlBattar}
    \newfontfamily\khatt[Script=Arabic, Scale=1.2]{AlBattar}
}{%
    \newfontfamily\arabicfontsf[
        Script=Arabic, Scale=1.2, Path=../fonts/,
        BoldFont=ae_AlBattar.ttf
    ]{ae_AlBattar.ttf}
    \newfontfamily\khatt[
        Script=Arabic, Scale=1.2, Path=../fonts/,
        BoldFont=ae_AlBattar.ttf
    ]{ae_AlBattar.ttf}
}

%------------------------------------------------
% Colors (same palette as the slides)
%------------------------------------------------
\definecolor{brandTeal}{HTML}{00B4D8}
\definecolor{brandOrange}{HTML}{FF7043}
\definecolor{brandDark}{HTML}{263238}
\definecolor{brandMuted}{HTML}{6C757D}
\colorlet{tealDeep}{brandTeal!70!black}

%------------------------------------------------
% Page setup
%------------------------------------------------
\renewcommand{\footrulewidth}{1pt}
\renewcommand{\headrulewidth}{1pt}
\rhead{\khatt المدرسة العليا للأساتذة ببوزريعة}
\lhead{\khatt ملحقة بومرداس}
\cfoot{\thepage}
\lfoot{\khatt الأستاذ: د. بن عبد الرزاق يوسف}
\parindent 0pt
\setlength{\headsep}{15pt}
\setlength{\headheight}{16pt}
\setlength{\parskip}{4pt}
\pagestyle{fancy}

\setlist[enumerate,1]{label=\arabic*), itemsep=4pt, topsep=3pt}
\setlist[enumerate,2]{label=\alph*), itemsep=2pt, topsep=2pt}
\setlist[itemize]{itemsep=2pt, topsep=2pt}

%------------------------------------------------
% Tables
%------------------------------------------------
\renewcommand{\arraystretch}{1.4}
\setlength{\arrayrulewidth}{0.6pt}
\arrayrulecolor{brandMuted!80}
% Header cell
\newcommand{\ras}[1]{\cellcolor{brandTeal!20}\textbf{#1}}
% Highlighted result cell
\newcommand{\natija}[1]{\cellcolor{brandOrange!15}\textbf{#1}}
% Empty cell to be filled by the student
\newcommand{\fragh}{\rule{0pt}{0.55cm}}
\newcolumntype{C}[1]{>{\centering\arraybackslash}m{#1}}
\newcolumntype{R}[1]{>{\raggedleft\arraybackslash}m{#1}}

%------------------------------------------------
% Boxes
%------------------------------------------------
\tcbset{labbox/.style={enhanced, breakable, arc=3pt, boxrule=1.2pt,
    coltitle=white, fonttitle=\khatt\bfseries,
    top=1.5mm, bottom=1.5mm, left=3mm, right=3mm,
    before skip=6pt, after skip=6pt}}
% Worked example
\newtcolorbox{mithal}[1]{labbox, colback=brandTeal!6, colframe=brandTeal,
    title={مثال محلول: #1}}
% Reminder / method
\newtcolorbox{tadhkir}[1][تذكير]{labbox, colback=brandTeal!3,
    colframe=tealDeep, title={#1}}
% Warning
\newtcolorbox{tanbih}[1][انتبه]{labbox, colback=brandOrange!8,
    colframe=brandOrange, title={#1}}
% Exercise
\newtcolorbox{tamrin}[1]{labbox, colback=white, colframe=brandDark!80,
    title={#1}}

% Section title
\newcommand{\qism}[1]{\section*{\khatt\underline{#1}}}

% Title block: \unwan{lab number}{title}{subtitle}
\newcommand{\unwan}[3]{%
    \begin{center}
        {\Large\khatt\textbf{المخبر التطبيقي رقم #1}}\\[0.15cm]
        {\Large\khatt\textbf{#2}}\\[0.1cm]
        {\khatt\textbf{#3}}
    \end{center}
    \rule{\textwidth}{1pt}
    \begin{center}
        \small المدة: ساعة ونصف \quad$\cdot$\quad الأدوات: آلة حاسبة، أداة ذكاء اصطناعي، مصدر موثوق
    \end{center}
    \rule{\textwidth}{1pt}}

% Percentage reminder used in every lab
\newcommand{\nisba}{%
    \begin{tadhkir}[تذكير: حساب النسبة المئوية]
        \begin{center}
            \textbf{النسبة المئوية = (الجزء $\div$ الكل) $\times$ 100}
        \end{center}
        مثال: 15 إجابة صحيحة من 20 سؤالًا:
        $\dfrac{15}{20}\times 100 = 75\%$.
    \end{tadhkir}}


\begin{document}

\unwan{1}{اكتشاف الذكاء الاصطناعي التوليدي}{أمثلة بالأرقام ــ التحقق ــ المقارنة بين الأدوات}

\qism{أهداف المخبر}
\begin{itemize}
    \item قياس صحة إجابات أداة ذكاء اصطناعي بالأرقام، وحساب نسبة الخطأ (الهلوسة).
    \item المقارنة بين طلب عام وطلب دقيق بمعايير قابلة للقياس.
    \item تصنيف المحتوى الذي تنتجه الأدوات التوليدية.
    \item التحقق من المعلومة في مصدر موثوق قبل استعمالها.
\end{itemize}

\begin{tanbih}[ملاحظة حول البيانات]
    الأرقام الواردة في هذا المخبر أمثلة تعليمية مبنية على تجارب في القسم،
    وليست تقييمًا رسميًا لأي أداة. أسماء الأدوات مُعوَّضة بالرموز (أ، ب، ج).
\end{tanbih}

\nisba

%================================================
\qism{الجزء الأول: أمثلة محلولة}
%================================================

\begin{mithal}{نسبة الإجابات الصحيحة}
طرح أستاذ \textbf{20 سؤالًا} في اللغة العربية على أداة ذكاء اصطناعي،
ثم صحّح الإجابات بالرجوع إلى كتب النحو والمعاجم والدواوين المحققة.

\begin{center}
\begin{tabular}{|R{4.6cm}|C{2.2cm}|C{2cm}|C{2cm}|C{2.6cm}|}
    \hline
    \ras{المجال} & \ras{عدد الأسئلة} & \ras{صحيحة} & \ras{خاطئة} & \ras{نسبة الصحة} \\ \hline
    النحو والإعراب              & 8  & 7  & 1 & $87.5\%$ \\ \hline
    الصرف                       & 4  & 3  & 1 & 75\% \\ \hline
    معاني الكلمات               & 3  & 3  & 0 & 100\% \\ \hline
    نسبة الأبيات إلى قائليها    & 5  & 2  & 3 & \natija{40\%} \\ \hline
    \textbf{المجموع}            & \textbf{20} & \textbf{15} & \textbf{5} & \natija{75\%} \\ \hline
\end{tabular}
\end{center}

\textbf{الحساب:}
\begin{itemize}
    \item نسبة الصحة الإجمالية: $\dfrac{15}{20}\times 100 = 75\%$،
          إذن نسبة الخطأ: $100 - 75 = 25\%$.
    \item في نسبة الأبيات: $\dfrac{2}{5}\times 100 = 40\%$ فقط، أي أن
          $\dfrac{3}{5}$ من الإجابات خاطئة ($60\%$).
\end{itemize}

\textbf{الاستنتاج:} الأداة جيدة في معاني الكلمات، لكنها تخطئ كثيرًا في نسبة الأبيات
إلى قائليها، وتكتب الخطأ بثقة. هذا هو \textbf{الهلوسة}. لذلك نتحقق دائمًا،
خاصة في المجال الذي تكثر فيه الأخطاء.
\end{mithal}

\begin{mithal}{طلب عام أم طلب دقيق؟}
كتب أستاذ طلبين في الموضوع نفسه، ثم قاس خصائص الإجابتين:

\textbf{الطلب العام:} «اشرح النعت.»\\
\textbf{الطلب الدقيق:} «اشرح النعت لتلميذ في السنة الأولى متوسط، بلغة بسيطة،
مع ثلاثة أمثلة من الحياة اليومية، في أقل من 150 كلمة.»

\begin{center}
\begin{tabular}{|R{5.5cm}|C{3cm}|C{3cm}|}
    \hline
    \ras{المعيار} & \ras{الطلب العام} & \ras{الطلب الدقيق} \\ \hline
    عدد كلمات الإجابة            & 420 & 140 \\ \hline
    عدد الأمثلة                  & 1   & 3   \\ \hline
    عدد المصطلحات الصعبة         & 9   & 1   \\ \hline
    تقييم الأستاذ (من 10)        & 4   & \natija{9} \\ \hline
\end{tabular}
\end{center}

\textbf{الحساب:}
\begin{itemize}
    \item الفرق في التقييم: $9 - 4 = 5$ نقاط.
    \item نقص عدد الكلمات: $420 - 140 = 280$ كلمة،
          أي $\dfrac{280}{420}\times 100 \approx 66.7\%$ أقصر.
\end{itemize}
\textbf{الاستنتاج:} الطلب الدقيق أعطى إجابة أقصر، وأنسب للمستوى، وفيها العدد المطلوب من الأمثلة.
\end{mithal}

%================================================
\qism{الجزء الثاني: تمارين تطبيقية}
%================================================

\begin{tamrin}{التمرين 1: المقارنة بين ثلاث أدوات}
طرح طالب \textbf{10 أسئلة} نفسها على ثلاث أدوات (أ، ب، ج)، وسجّل عدد الإجابات الصحيحة:

\begin{center}
\begin{tabular}{|R{3.4cm}|C{2.2cm}|C{2cm}|C{2cm}|C{2cm}|}
    \hline
    \ras{المجال} & \ras{عدد الأسئلة} & \ras{الأداة أ} & \ras{الأداة ب} & \ras{الأداة ج} \\ \hline
    النحو              & 4  & 4 & 3 & 2 \\ \hline
    البلاغة            & 3  & 2 & 3 & 2 \\ \hline
    نسبة الأبيات       & 3  & 1 & 2 & 0 \\ \hline
    \textbf{المجموع}   & \textbf{10} & \fragh & \fragh & \fragh \\ \hline
    \textbf{نسبة الصحة} & --  & \fragh & \fragh & \fragh \\ \hline
\end{tabular}
\end{center}

\begin{enumerate}
    \item أكمل السطرين الأخيرين من الجدول.
    \item رتّب الأدوات الثلاث من الأفضل إلى الأضعف.
    \item ما المجال الذي أخطأت فيه الأدوات الثلاث أكثر؟ احسب نسبة الصحة فيه
          لكل الأدوات معًا (مجموع الإجابات الصحيحة $\div$ 9).
    \item هل تعني النتيجة الأفضل أننا لا نحتاج إلى التحقق من إجابات هذه الأداة؟ علّل.
\end{enumerate}
\end{tamrin}

\begin{tamrin}{التمرين 2: تحقق من المعلومة}
أعطت أداة ذكاء اصطناعي المعلومات الآتية. تحقق من كل معلومة في مصدر موثوق
(ديوان، معجم، كتاب نحو)، ثم أكمل الجدول:

\begin{center}
\begin{tabular}{|C{0.6cm}|R{6.6cm}|C{2cm}|R{3.6cm}|}
    \hline
    \ras{ر} & \ras{ما قالته الأداة} & \ras{صحيح / خطأ} & \ras{المعلومة الصحيحة} \\ \hline
    1 & «قِفا نبكِ من ذكرى حبيبٍ ومنزلِ» لعنترة بن شداد. & \fragh & \\ \hline
    2 & «الخيلُ والليلُ والبيداءُ تعرفني» للمتنبي. & \fragh & \\ \hline
    3 & «إذا الشعبُ يومًا أراد الحياة» لأحمد شوقي. & \fragh & \\ \hline
    4 & مؤلف معجم «لسان العرب» هو ابن منظور. & \fragh & \\ \hline
    5 & كتاب «الكتاب» في النحو لسيبويه. & \fragh & \\ \hline
    6 & «الألفية» في النحو منظومة لابن هشام. & \fragh & \\ \hline
\end{tabular}
\end{center}

\begin{enumerate}
    \item كم معلومة خاطئة وجدتَ؟ احسب نسبة الهلوسة في هذه الإجابة.
    \item اذكر المصدر الذي اعتمدت عليه في كل تصحيح.
    \item لو نقل أستاذ هذه المعلومات إلى تلاميذه دون تحقق، كم معلومة خاطئة سيتعلمون؟
\end{enumerate}
\end{tamrin}

\begin{tamrin}{التمرين 3: ماذا تنتج الأداة؟}
صنّف كل طلب حسب نوع المحتوى: \textbf{نص}، \textbf{صورة}، \textbf{صوت}،
\textbf{فيديو}، \textbf{جدول أو مخطط}.

\begin{center}
\begin{tabular}{|C{0.6cm}|R{9.5cm}|C{3.5cm}|}
    \hline
    \ras{ر} & \ras{الطلب} & \ras{نوع المحتوى} \\ \hline
    1  & قصة قصيرة لتلاميذ السنة الأولى متوسط. & \fragh \\ \hline
    2  & رسم يوضح مشهدًا من نص القراءة. & \fragh \\ \hline
    3  & قراءة مسجلة لقصيدة مع ضبط الحركات. & \fragh \\ \hline
    4  & جدول يلخص أنواع الخبر مع أمثلة. & \fragh \\ \hline
    5  & خمسة أسئلة لفهم نص القراءة. & \fragh \\ \hline
    6  & مقطع قصير متحرك يشرح كتابة الهمزة. & \fragh \\ \hline
    7  & ملخص لدرس الجملة الاسمية. & \fragh \\ \hline
    8  & مخطط يبيّن أقسام الكلمة: اسم، فعل، حرف. & \fragh \\ \hline
    9  & صورة غلاف لمجلة القسم. & \fragh \\ \hline
    10 & تسجيل للنطق الصحيح لمخارج الحروف. & \fragh \\ \hline
\end{tabular}
\end{center}

ثم أكمل جدول التكرارات:
\begin{center}
\begin{tabular}{|R{1.8cm}|C{1.5cm}|C{1.5cm}|C{1.5cm}|C{1.5cm}|C{2.4cm}|C{1.7cm}|}
    \hline
    \ras{النوع} & \ras{نص} & \ras{صورة} & \ras{صوت} & \ras{فيديو} & \ras{جدول أو مخطط} & \ras{المجموع} \\ \hline
    العدد  & \fragh & & & & & 10 \\ \hline
    النسبة & \fragh & & & & & 100\% \\ \hline
\end{tabular}
\end{center}
\end{tamrin}

\begin{tamrin}{التمرين 4: هل تربح الوقت حقًا؟}
يحضّر أستاذ \textbf{10 تمارين} في النحو:
\begin{itemize}
    \item \textbf{دون أداة:} 60 دقيقة.
    \item \textbf{مع أداة:} كتابة الطلب 5 دقائق، وتوليد الإجابة دقيقة واحدة،
          وقراءة التمارين وتصحيحها 19 دقيقة.
\end{itemize}
\begin{enumerate}
    \item احسب الزمن الكلي مع الأداة.
    \item كم دقيقة ربح الأستاذ؟ وما نسبة هذا الربح من 60 دقيقة؟
    \item وجد الأستاذ أثناء المراجعة 3 تمارين خاطئة من 10. لو لم يراجع، وحضّر 4 سلاسل
          في الأسبوع بالطريقة نفسها، كم تمرينًا خاطئًا سيصل إلى التلاميذ في الأسبوع؟
    \item ماذا تستنتج عن دور المراجعة؟
\end{enumerate}
\end{tamrin}

\begin{tamrin}{التمرين 5: تجربة في القسم}
اختر قاعدة نحوية (مثل: الحال، التمييز، المفعول المطلق). اكتب لها طلبًا عامًا
وطلبًا دقيقًا، ونفّذهما في الأداة نفسها، ثم أكمل الجدول:

\begin{center}
\begin{tabular}{|R{5.5cm}|C{3cm}|C{3cm}|}
    \hline
    \ras{المعيار} & \ras{الطلب العام} & \ras{الطلب الدقيق} \\ \hline
    عدد كلمات الإجابة            & \fragh & \\ \hline
    عدد الأمثلة                  & \fragh & \\ \hline
    عدد الأخطاء التي وجدتها       & \fragh & \\ \hline
    تقييمك (من 10)               & \fragh & \\ \hline
\end{tabular}
\end{center}
احسب الفرق في التقييم، واكتب جملتين تشرح فيهما سبب الفرق.
\end{tamrin}

%================================================
\newpage
\qism{الحلول المختصرة}
%================================================

\textbf{التمرين 1:}
المجموع: أ $= 7$، ب $= 8$، ج $= 4$. النسب: أ $70\%$، ب $80\%$، ج $40\%$.
الترتيب: ب ثم أ ثم ج.
المجال الأضعف: نسبة الأبيات، مجموع الصحيح $1+2+0=3$ من 9، أي $\approx 33.3\%$.
لا، حتى الأداة الأفضل أخطأت في سؤالين من 10، فالتحقق ضروري دائمًا.

\textbf{التمرين 2:}
1) خطأ: لامرئ القيس. 2) صحيح. 3) خطأ: لأبي القاسم الشابي. 4) صحيح. 5) صحيح.
6) خطأ: لابن مالك.
عدد الأخطاء 3 من 6، نسبة الهلوسة $\dfrac{3}{6}\times 100 = 50\%$.
سيتعلم التلاميذ 3 معلومات خاطئة.

\textbf{التمرين 3:}
نص: 1، 5، 7 \quad صورة: 2، 9 \quad صوت: 3، 10 \quad فيديو: 6 \quad جدول أو مخطط: 4، 8.\\
العدد: 3، 2، 2، 1، 2 \quad النسبة: $30\%$، $20\%$، $20\%$، $10\%$، $20\%$.

\textbf{التمرين 4:}
الزمن مع الأداة: $5 + 1 + 19 = 25$ دقيقة.
الربح: $60 - 25 = 35$ دقيقة، أي $\dfrac{35}{60}\times 100 \approx 58.3\%$.
الأخطاء دون مراجعة: $3 \times 4 = 12$ تمرينًا خاطئًا في الأسبوع.
المراجعة تأخذ أطول جزء من الوقت (19 دقيقة)، لكنها هي التي تحمي التلاميذ من الأخطاء.

\textbf{التمرين 5:} إجابة شخصية؛ يُقبل كل جدول مملوء بقياسات حقيقية مع تعليل منطقي.

\end{document}
